\chapter{Supplementary Regression Experiments}
\label{app:supplementary-regression-experiments}

This appendix keeps the regression-track details needed to audit the main
claims without reproducing the full run logs. The detailed source artifacts and
analysis scripts remain in the repository; the compiled thesis keeps only the
protocol and summary information that supports Chapter~\ref{chap:regression-curriculum}.

\section{Hyperparameter sweep baseline}
\label{app:sweep-pack-v1}

Before curriculum comparisons, each benchmark function received its own Optuna
search over MLP, Fourier-feature, and SIREN architectures. This was necessary to
avoid comparing curriculum runs against under-tuned direct baselines. The sweep
selected different architecture families for different functions, which justified
using per-function best configurations in the fixed continuation benchmark.

\begin{table}[h]
\centering
\small
\begin{tabular}{lrrl}
\toprule
Function & Best validation loss & Parameters & Best architecture \\
\midrule
\texttt{ackley} & 0.069 & 1,185 & MLP \\
\texttt{bukin} & 57.9 & 469 & MLP \\
\texttt{eggholder} & 1.22e4 & 4,545 & Fourier \\
\texttt{franke} & 1.77e-4 & 1,265 & MLP \\
\texttt{levy} & 0.376 & 3,745 & MLP \\
\texttt{peaks} & 0.010 & 737 & SIREN \\
\bottomrule
\end{tabular}
\caption{Representative best configurations from the baseline sweep. The exact
numbers are used only to establish that the direct baselines were tuned rather
than manually chosen defaults.}
\label{tab:app-regression-sweep-summary}
\end{table}

\section{v3 ratio-constrained dataset setting}
\label{app:sweep-pack-v3-setup}

The stricter v3 setting used 20,000 generated samples per function, split into
10,000 train and 10,000 test samples, with an internal 80/20 train--validation
split inside the training file. Training labels included controlled Gaussian
noise, and model searches enforced a parameter-to-sample ratio constraint. This
setting made the regression experiments less vulnerable to overlarge models
memorizing small synthetic datasets.

The active training pipelines logged validation loss, final test loss where
available, gradient and stability diagnostics, configuration files, and dataset
references to MLflow. These artifacts are the source for the regression figures
and tables in the main text.

\begin{table}[h]
\centering
\scriptsize
\setlength{\tabcolsep}{3pt}
\begin{tabular}{@{}p{0.14\textwidth}p{0.23\textwidth}p{0.25\textwidth}p{0.28\textwidth}@{}}
\toprule
Version & Purpose & Data setting & Main use in the thesis \\
\midrule
v1 sweep pack & Tune baselines before curriculum comparison & Function-specific sample counts, 5--6 smoothing levels, no ratio constraint & Supplies per-function architectures for the fixed continuation benchmark. \\
v1 curriculum benchmark & Compare direct and fixed continuation training & Same v1 data, two budget regimes, three seeds & Main evidence for mixed final-loss effects and runtime cost. \\
v2 ratio sweeps & Check capacity control & 20,000 samples, 3 smoothing levels, 2\% label noise, 10--20 train samples per parameter & Tests whether baseline conclusions survive a parameter-to-data constraint. \\
v3 ratio sweeps & Stress-test the stricter setting & Same as v2 but 3\% label noise & Source of the MLflow ratio-constrained diagnostics and architecture distribution checks. \\
\bottomrule
\end{tabular}
\caption{Regression experiment versions. The table separates baseline tuning,
fixed curriculum comparison, and later ratio-constrained checks so that v1, v2,
and v3 are not read as a single interchangeable experiment.}
\label{tab:app-regression-version-summary}
\end{table}

\section{Fixed continuation benchmark}
\label{app:curr-vs-single-v1}

The fixed benchmark compared direct hard-target training with a fixed
continuation schedule under two budgets: equal epochs and equal wall-clock time.
It used eight functions and three random seeds for each method/budget
combination. The main text reports function-level loss differences and the
quality/runtime trade-off. The essential protocol constraints were:

\begin{itemize}
    \item same train/test data and hard-validation metric per function,
    \item same tuned architecture family and optimizer configuration per
          function,
    \item only the training strategy changed between direct and curriculum runs,
    \item runtime was recorded separately from final quality.
\end{itemize}

This design is why the negative runtime result is meaningful: curriculum was not
slower because it used a larger model, but because the staged training procedure
itself added cost.

\section{Meta-weighting benchmark}
\label{app:meta-weighting-eggholder-u1}

The reduced meta-weighting benchmark used \texttt{eggholder} as the difficult
representative target and tested one-step unrolled meta-updates over
continuation losses. The grid covered four seeds, three noise levels, several
model and meta learning rates, and \(1,2,3,5\) losses. All 288 runs completed.

The main text reports that all multi-loss settings were negative on average
against the one-loss hard-target baseline. The supplementary point is that the
method was operationally feasible and produced interpretable schedules, but the
outer-loop signal was too noisy to make the result robust across seeds and noise
levels. This is why meta-weighting is treated as future work rather than a
successful replacement for fixed schedules.
